% ============================================================
%  Part X — Comparison
% ============================================================
\gpart{X}{Comparison}

% ------------------------------------------------------------
\gsec{Comparative and superlative}

\gintro{German compares by adding endings rather than by using \emph{more} and \emph{most}, and it does so regardless of how long the adjective is.}

\begin{gtbl}{colspec={X[0.9,l] X[1,l] X[1.1,l]}}
  Degree & Formation & Example \\
  Positive    & the plain adjective & schnell \\
  Comparative & + \en{-er}          & schnell\textbf{er} \\
  Superlative & am + \en{-sten}     & \textbf{am} schnell\textbf{sten} \\
  Superlative & before a noun: \en{-ste} + ending & der schnell\textbf{ste} Wagen \\
\end{gtbl}

\begin{gnote}
German adds endings where English reaches for \emph{more} and \emph{most}. Even
long adjectives follow the rule: \textit{interessanter},
\textit{am interessantesten}. There is no \textit{mehr interessant}.
\end{gnote}

\tcap{Umlaut in one-syllable adjectives}

\begin{gtbl}{colspec={X[1,l] X[1,l] X[1,l]}}
  Positive & Comparative & Superlative \\
  alt   & älter   & am ältesten \\
  jung  & jünger  & am jüngsten \\
  lang  & länger  & am längsten \\
  kurz  & kürzer  & am kürzesten \\
  stark & stärker & am stärksten \\
  warm  & wärmer  & am wärmsten \\
  kalt  & kälter  & am kältesten \\
  groß  & größer  & am größten \\
\end{gtbl}

\begin{gnote}
Insert an extra \textbf{e} before \textit{-sten} when the adjective ends in
\textit{-d, -t, -s, -ß, -z, -sch}: \textit{am kält\textbf{e}sten},
\textit{am heiß\textbf{e}sten}. \textit{groß} is the exception —
\textit{am größten}.
\end{gnote}

% ------------------------------------------------------------
\gsec{Irregular comparisons}

\gintro{A handful of very common words compare irregularly. They are the same words that are irregular in English.}

\begin{gtbl}{colspec={X[0.9,l] X[0.9,l] X[1.1,l] X[0.8,l]}}
  Positive & Comparative & Superlative & Meaning \\
  gut   & besser & am besten   & good \\
  viel  & mehr   & am meisten  & much, many \\
  gern  & lieber & am liebsten & gladly \\
  hoch  & höher  & am höchsten & high \\
  nah   & näher  & am nächsten & near \\
  groß  & größer & am größten  & big \\
  bald  & eher   & am ehesten  & soon \\
  wenig & weniger / minder & am wenigsten & little \\
\end{gtbl}

\begin{gnote}
\textit{gern} $\rightarrow$ \textit{lieber} $\rightarrow$ \textit{am liebsten}
is how German expresses preference: \textit{Ich trinke \textbf{gern} Tee, aber
\textbf{lieber} Kaffee, und \textbf{am liebsten} Bier.}
\end{gnote}

% ------------------------------------------------------------
\gsec{Comparison patterns}

\gintro{Beyond the plain comparative there are a few fixed frames for expressing equality, gradual change and proportion.}

\begin{gtbl}{colspec={X[0.95,l] X[1.05,l] X[1,l]}}
  Pattern & Meaning & Example \\
  (genau)so \dots\ wie & as \dots\ as       & Er ist \textbf{so groß wie} ich. \\
  nicht so \dots\ wie  & not as \dots\ as   & Sie ist \textbf{nicht so alt wie} du. \\
  \dots-er \textbf{als} & \dots-er than     & Er ist \textbf{größer als} ich. \\
  immer \dots-er       & \dots\ and \dots-er & Es wird \textbf{immer kälter}. \\
  je \dots, desto \dots & the \dots, the \dots & \textbf{Je mehr}, \textbf{desto besser}. \\
\end{gtbl}

\begin{gnote}
\textbf{als} for a difference, \textbf{wie} for equality. \textit{größer als},
never \textit{größer wie} — that is a very common native slip, but it is not
standard.
\end{gnote}

\tcap{je \dots\ desto}

\begin{flattbl}{colspec={X[1,l] X[1.6,l]}}
  Structure & \textit{je} + comparative, verb LAST; \textit{desto} + comparative, verb SECOND \\
  Example   & \textbf{Je} früher du \textbf{kommst}, \textbf{desto} besser \textbf{ist} es. \\
  Example   & \textbf{Je} mehr ich \textbf{lerne}, \textbf{desto} weniger \textbf{weiß} ich. \\
\end{flattbl}

\tcap{Superlative before a noun takes adjective endings}

\begin{gendertbl}{rows={font=\footnotesize}}
  der schnellste & maskulin & feminin & neutrum & Plural \\
  Nominativ & der schnellste & die schnellste & das schnellste & die schnellsten \\
  Akkusativ & den schnellsten & die schnellste & das schnellste & die schnellsten \\
  Dativ     & dem schnellsten & der schnellsten & dem schnellsten & den schnellsten \\
\end{gendertbl}

\begin{gnote}
Use \textit{am schnellsten} when the superlative stands alone after the verb,
and \textit{der/die/das schnellste} when it sits in front of a noun or replaces
one: \textit{Er läuft \textbf{am schnellsten}} vs \textit{Er ist \textbf{der
schnellste} Läufer}.
\end{gnote}
